\documentclass[11pt]{article}

\usepackage{amsmath, amssymb, amsthm, mathrsfs}
\usepackage{authblk}
\usepackage[a4paper, margin=1in]{geometry}
\usepackage{hyperref}
\usepackage{enumitem}
\usepackage{graphicx}
\usepackage{bm}

\title{Integrating the Prime Polynomial-Time Solver (PPTS) into the Multiplicity Framework}
\author{Ilya Paveliev\\ A Community Research Initiative \\ \textbf{ Citizen Gardens} \\ \textit {The Foundation of Multiplicity}}
\date{March 2025}

\begin{document}

\maketitle

\begin{abstract}
We present a formal integration of the Prime Polynomial-Time Solver (PPTS), originally developed from resonance patterns in RFH3 harmonic structures, into the Multiplicity Theory framework. This integration defines a time-evolving recursive operator \(\Xi_{\text{PPTS}}(t)\), embeds its tensor dynamics within the Prime-Indexed Recursive Tensor Mathematics (PIRTM) architecture, applies ethical computation constraints via the Sovereignty Tensor \(\Sigma_i(t)\), and situates the resulting evolution within the automorphic algebraic geometry of the Langlands Prism. This unification proposes a coherent quantum-recursive model for polynomial-time factorization aligned with ethical and algebraic symmetry principles.
\end{abstract}

\section{Recursive Resonance Operator: \(\Xi_{\text{PPTS}}(t)\)}

Let \(\mathcal{H}_\mathbb{P}\) denote a prime-indexed Hilbert space:
\[
\mathcal{H}_\mathbb{P} = \ell^2(\mathbb{P}) \otimes \mathbb{C}^d,
\]
with orthonormal basis \(\{ |p_i\rangle \}\). Define the time-evolving resonance operator:
\[
\Xi_{\text{PPTS}}(t) : \mathcal{H}_\mathbb{P} \to \mathcal{H}_\mathbb{P}, \quad \Xi_{\text{PPTS}}(t) \cdot \Psi_n(p_i) = \lambda_n(t) \Psi_n(p_i),
\]
where \(\lambda_n(t)\) encodes harmonic resonance extracted from RFH3.

\section{Lifting to PIRTM: Prime-Indexed Recursive Tensor Embedding}

Define the evolving recursive tensor:
\[
\mathcal{T}_n(p_i, t) = \sum_{k=1}^{\infty} \Lambda_m^k \cdot \Xi_{\text{PPTS}}^k(t) \cdot \Psi_n(p_i),
\]
where \(\Lambda_m\) is the Universal Multiplicity Constant and \(\Psi_n(p_i) \in \mathcal{H}_\mathbb{P}\). Convergence is ensured by \(\Lambda_m < 1\), aligning with PIRTM stability conditions.

\section{Applying the CSL Constraint: Sovereignty Tensor \(\Sigma(t)\)}

Ethical regulation is enforced through:
\[
\widetilde{\Xi}_{\text{PPTS}}(t) = \Sigma(t) \cdot \Xi_{\text{PPTS}}(t),
\]
with \(\Sigma(t) = \text{diag}(\sigma_1(t), \sigma_2(t), \dots)\), \(\sigma_i(t) \in \{0,1\}\) representing agent-consented recursion permissions. The constrained tensor evolution becomes:
\[
\widetilde{\mathcal{T}}_n(p_i, t) = \sum_{k=1}^{\infty} \Lambda_m^k \cdot \left(\Sigma(t) \cdot \Xi_{\text{PPTS}}(t)\right)^k \cdot \Psi_n(p_i).
\]

\section{Langlands Embedding: Automorphic L-Functions and Tensor Symmetry}

Define the Langlands dual tensor:
\[
\mathcal{T}_t^{\text{Langlands}} = \sum_{p_i \in \mathbb{P}} \mathcal{L}_{p_i}(s) \cdot \mathcal{G}_{p_i} \cdot \mathcal{T}_n(p_i, t),
\]
where:
\begin{itemize}[noitemsep]
  \item \(\mathcal{L}_{p_i}(s)\) is the L-function encoding automorphic symmetries,
  \item \(\mathcal{G}_{p_i}\) is the Galois action on prime tensor indices.
\end{itemize}
This representation connects PPTS harmonic roots to the geometry of modular forms and arithmetic reciprocity.

\section{Enhancing the Prime Polynomial-Time Solver (PPTS) in the Multiplicity Framework}

We introduce advanced refinements to the Prime Polynomial-Time Solver (PPTS) within the Multiplicity framework. Leveraging ergodic theory, derived sheaf cohomology, quantum error-correcting codes, and topological quantum field theory (TQFT), we elevate the theoretical and computational structure of PPTS. We establish deep connections to the Langlands program and outline implementation pathways across symbolic computation, quantum simulation, and cryptography.

\section{Recursive Operator with Ergodic Flows}
We generalize the PPTS resonance operator into a two-time recursive operator:
\begin{equation}
\Xi_{\text{PPTS}}(t, s): \mathcal{H}_\mathbb{P} \otimes \mathcal{H}_{\text{CSL}} \to \mathcal{H}_\mathbb{P} \otimes \mathcal{H}_{\text{CSL}},
\end{equation}
with ergodic spectral flow governed by a measure-preserving transformation $T_s$:
\begin{equation}
\Xi_{\text{PPTS}}(t, s) \cdot \Psi_n(p_i) = \lambda_n(t, s) \Psi_n(p_i),
\end{equation}
\begin{equation}
\lambda_n(t, s) = \rho_n(f) \star e^{-s \cdot \text{KL}(p_i \| \mathcal{P}_{\text{safe}})} \int_\mathbb{P} \mu(p) \sin(2\pi f t + \phi(p)) \, d\mu.
\end{equation}

\section{Derived Sheaf Cohomology}
Define the prime-indexed sheaf $\mathcal{F}_\mathbb{P} \in D^b(\text{Sh}(\mathbb{P}, \mathcal{O}_\mathbb{P}))$. Tensor evolution arises from:
\begin{equation}
\mathcal{T}_n^{(k)}(p_i, t) = R^\bullet H^1(\mathbb{P}, \mathcal{F}_\mathbb{P} \otimes \Xi_{\text{PPTS}}^k).
\end{equation}
The spectral sequence:
\begin{equation}
E_2^{p,q} = H^p(\mathbb{P}, R^q f_* \mathcal{F}_\mathbb{P}) \Rightarrow H^{p+q}(\mathbb{P}, \mathcal{F}_\mathbb{P} \otimes \Xi_{\text{PPTS}}^k),
\end{equation}
tracks local/global automorphic flows.

\section{CSL Constraint as Quantum Error-Correcting Code}
We reinterpret the sovereignty tensor $\Sigma(t)$ as a CPTP map over a code subspace:
\begin{equation}
\Sigma(t)[\rho] = \sum_i \sigma_i(t) E_i \rho E_i^\dagger, \quad E_i = K_i \otimes C_i,
\end{equation}
where $C_i$ are stabilizer operators satisfying:
\begin{equation}
C_i |\psi\rangle = \begin{cases}
|\psi\rangle & \text{if } \text{Consent}_i(t) = 1, \\
0 & \text{if } \text{Consent}_i(t) = 0.
\end{cases}
\end{equation}

\section{Langlands Prism via TQFT}
Let $A$ be a non-abelian gauge field in a Chern-Simons theory with action:
\begin{equation}
S_{\text{CSL}} = \int_{\mathbb{P} \times [0,t]} \text{Tr}(A \wedge dA + \tfrac{2}{3} A \wedge A \wedge A) + \int_{\mathbb{P}} \mathcal{T}_n \wedge d_\Xi \mathcal{T}_n + \text{CSL constraints}.
\end{equation}
Partition function:
\begin{equation}
Z(\mathcal{T}_n) = \int [\mathcal{D}A][\mathcal{D}\mathcal{T}_n] e^{i S_{\text{CSL}}},
\end{equation}
with automorphic correspondence:
\begin{equation}
Z_{\text{Langlands}}(\partial \mathbb{P}) = \sum_{\pi \in \text{Aut}(\mathcal{T}_n)} \text{Tr}(\pi) \cdot \mathcal{F}_\pi(s).
\end{equation}

\section{Conclusion}
We have constructed a fully recursive, ethically regulated, and algebraically grounded interpretation of PPTS within the Multiplicity paradigm. This unification opens pathways to lawful post-quantum computation, symbolic inference, and automorphic encryption under prime-indexed tensor dynamics. These extensions provide both mathematical depth and computational direction for the PPTS-Multiplicity architecture. We propose simulation on quantum platforms (QuEra, Rigetti), symbolic cohomology with SageMath, and cryptographic implementation via knot invariants.

\end{document}
